\documentclass[letterpaper]{article}
\usepackage[submission]{aaai2027}
\usepackage[hyphens]{url}
\usepackage{graphicx}
\urlstyle{rm}
\def\UrlFont{\rm}
\usepackage{natbib}
\usepackage{caption}
\frenchspacing
\usepackage{amsmath,amssymb,amsfonts}
\usepackage{booktabs}
\pdfinfo{/TemplateVersion (2027.1)}
\setcounter{secnumdepth}{2}
\newcommand{\name}{ScaleDown}
\newcommand{\pending}[1]{\textbf{[#1: pending]}}
\newcommand{\slot}[2]{\noindent\fbox{\parbox{\dimexpr\linewidth-2\fboxsep-2\fboxrule\relax}{\small\textbf{#1 --- experiment placeholder.} #2}}}
\newcommand{\figslot}[2]{\fbox{\parbox[c][#1][c]{0.94\linewidth}{\centering\small\textbf{RESULT FIGURE PLACEHOLDER}\\[4pt]#2}}}
\newtheorem{proposition}{Proposition}

\title{\name{}: Learning Queue-Aware Resource Requests\\
for Carbon-Efficient Training under Completion Budgets}
\author{Anonymous Submission}
\affiliations{}
\begin{document}
\maketitle
\begin{abstract}
An ordinary user of a shared GPU cluster controls resource requests, while the
batch scheduler controls when they run. This separation makes carbon-aware
training a sequential decision problem: node count changes training efficiency,
queue wait, and the carbon intensity encountered during execution. We present
\name{}, a controller that selects the node count of each checkpointed training
chunk under a slider-selected completion budget. Its policy combines remaining
work and time with public queue history and a causal carbon-intensity forecast
sequence, without requiring an admission-time predictor. An undiscounted, constrained actor--critic
objective accounts for variable chunk durations and deadline misses. To handle
unmeasured node power, we factor carbon accounting into execution
exposure and an explicit power scenario. Within a three-factor power model,
two endpoint evaluations characterize the full sensitivity interval, enabling
one policy to optimize across that interval. The evaluation compares against
all fixed scales, a fixed-policy mixture, and a queue-aware planning baseline
on two trace-derived demand scenarios with published AMSP throughput profiles
\pending{E1: scenario and input validation}.
\pending{E2: baseline-relative carbon and deadline results};
\pending{E3--E4: mechanism and robustness results}.
The resulting claims concern modeled job-attributed operational carbon;
facility emissions and production deployment are outside the evaluation.
\end{abstract}

\section{Introduction}

The operational carbon associated with training depends on electricity
consumption and when that electricity is used~\cite{patterson2021carbon}.
Carbon-aware systems exploit temporal, spatial, and allocation
flexibility~\cite{hanafy2023carbonscaler,sukprasert2024limitations}. On a shared
batch cluster, however, a user controls a \emph{request}, while Slurm determines
admission, placement, and start time. A resource choice may take effect hours
later, when the carbon intensity differs from that observed at submission.

Node count couples three consequences. Fewer nodes can reduce node-hours under
sublinear strong scaling, but lengthen execution and increase the number of
queue admissions needed by a long run. A smaller request may start sooner,
yet overlap a less favorable carbon interval. Near completion, a faster final
chunk may avoid another queue wait. The useful decision therefore depends on
remaining work, available time, and the execution window of each request.

The central question is \emph{when adapting requests provides value beyond a
good fixed scale}. With constant carbon intensity, no waiting, and no chunk
overhead, we show that mixed-scale execution lies in the convex hull of
fixed-scale time--carbon points. Scaling inefficiency alone offers no gain
below that frontier. Under deadlines, a preplanned within-run mix can still
outperform a per-run randomized fixed policy without new queue observations.
We therefore retain both fixed mixtures and a plan chosen once at submission
when testing the value of online adaptation.

\name{} exposes a completion-budget slider and makes repeated node-count decisions without
modifying Slurm. Its main actor--critic policy uses visible queue history,
remaining work and time, action execution descriptors, and a causal CI forecast
sequence. It requires no explicit wait predictor: continuation costs are
learned from realized replay outcomes, and each checkpoint reveals how much
budget remains. Hidden admissions still limit predictability and any deadline
claim. We compare against explicit rollout planning and a matched learned policy
that commits its full sequence before submission, separating dynamic feedback
from the benefit of a preplanned mix.

Complete-node power is unavailable, so we separate execution evidence from
carbon estimates. The power model has three factors: a normalization
reference, an unknown workload coefficient, and a scale coefficient. For a fixed
execution log, carbon is affine in the scale-model parameter. Two endpoint
comparisons thus characterize robustness within the declared interval,
without determining absolute power or repeating scheduler simulations.
We use this property to train against both endpoint objectives and report
where comparative conclusions change.

Our evaluation uses published AMSP performance profiles and two historical
arrival streams to define controlled contention scenarios. Hardware labels do
not determine admission once resource demands and scheduling rules are fixed;
this abstraction does not predict how old jobs would run on new hardware.

The paper makes three contributions:
\begin{itemize}
\item A completion-budget formulation and fixed-scale null case that separate
scaling efficiency, preplanned mixing, and online adaptive value.
\item A user-level controller combining realized budget feedback, public queue
history, and undiscounted constrained learning without an admission predictor.
\item An exposure-based sensitivity analysis and a matched evaluation against
fixed, plan-once, and replanning baselines.
\end{itemize}

\section{Related Work}

\paragraph{Carbon-aware elasticity.}
CarbonScaler varies a batch workload's cloud allocation with carbon intensity
and uses a greedy marginal-resource allocation algorithm
~\cite{hanafy2023carbonscaler}. It already establishes resource elasticity as a
carbon-control mechanism. \name{} addresses the additional delay between
requesting resources and obtaining them, as well as repeated admission for
checkpointed chunks. A direct transplant of cloud autoscaling would change
the available controls; our controls retain repeated admissions and charge every queue wait. Studies of temporal and spatial shifting
also show that flexibility and overhead constrain attainable carbon savings
~\cite{sukprasert2024limitations}. We report the scope of the opportunity rather
than assume that lower-carbon timing is always available.

\paragraph{User-side provisioning.}
QBETS estimates wait bounds from historical time series~\cite{nurmi2007qbets};
our main policy learns completed-run costs without an admission estimator.
Mirage uses learning to provision follow-up resources for low-interruption
services on batch GPU clusters~\cite{ding2023mirage}. We use its trace-replay
infrastructure, but optimize the node count of a ready training chunk under a
completion budget. Scheduler-level methods such as DRAS act on the scheduling
policy itself~\cite{fan2021deep}. Our policy has no such authority and cannot
observe hidden priority or reservation state.

\paragraph{Constrained learning and accounting uncertainty.}
Constrained policy optimization separates a resource objective from a service
constraint~\cite{achiam2017cpo}; PPO supplies a practical actor--critic optimizer
~\cite{schulman2017ppo}. We use these established ideas, without claiming a new
general-purpose RL optimizer or a constraint-satisfaction theorem. The method
contribution is the request-conditioned decision representation and its
coupling to a tractable power-scenario objective. Unlike operational-carbon
measurement, the latter assesses consequences conditional on a model. An
algebraic sensitivity certificate does not validate the model's physical
parameters.

\section{Problem and Adaptive Opportunity}
\label{sec:problem}

\paragraph{Control and work.}
A target run consists of $U_\star$ prescribed optimizer updates, executed as
sequential checkpointed chunks. At boundary $k$, time is $t_k$ and remaining
work is $U_k$. The action is a feasible node count
$n_k\in\mathcal A=\{4,16,64,128\}$. The next chunk is submitted immediately;
there is no intentional deferral, concurrent reservation, or resizing of a
running chunk. Scale choices may rise or fall between allocations; ``down
means using fewer than the maximum supported nodes, not irreversible shrinking. We use $H=48$ hours as the study's maximum requested wall time,
subject to the chosen partition's limits; it is not a universal Slurm limit.

For sharded or replicated data parallelism without tensor or pipeline
parallelism, let $g$ be GPUs per node ($g=8$ for the published AMSP profiles),
$b$ the per-GPU micro-batch, and $a_n$ accumulation. Each
action satisfies
\begin{equation}
 B_\star=ngba_n,\qquad a_n=B_\star/(ngb)\in\mathbb N.
 \label{eq:batch}
\end{equation}
The same data order, sequence length, optimizer-update target, and update-indexed
learning-rate schedule apply across actions. Accumulation implements the fixed
batch~\cite{shoeybi2019megatron}; it does not guarantee bitwise equality or equal
final quality. A configuration that fails the batch, memory, or partition
constraints is masked for every method.

Published steady-state throughput determines $q_n$, the updates per hour.
Initialization/restart $r_n$ and complete-checkpoint time $h_n$ are declared
scenarios, not measurements supplied by the throughput source. Define
\begin{align}
 v_{k,n}&=\min\{U_k,\lfloor q_n(H-r_n-h_n)\rfloor\},\nonumber\\
 d_{k,n}&=r_n+v_{k,n}/q_n+h_n,\label{eq:chunk}\\
 U_{k+1}&=U_k-v_{k,n_k},\quad
 t_{k+1}=t_k+w_k+d_{k,n_k}.\nonumber
\end{align}
Only positive-progress actions are allowed. The request's wall time is
$d_{k,n}$ rounded up to the scheduler's time granularity, capped at $H$; the
job releases resources on completion. Thus the final partial chunk has a
shorter request, which can change admission. The assumed restart contract includes model,
optimizer, learning-rate state, global update index, random-state metadata,
and the next global sample position; this study does not implement or validate
AMSP state migration. The final output save is included in
$h_n$. End-to-end turnaround time is $T=t_K-t_0$, including every queue wait
and allocation interval.

\paragraph{A fixed-scale null case.}
The following observation defines what scaling profiles alone can establish.

\begin{proposition}[No improvement below the fixed hull]
\label{prop:hull}
Suppose update rates and per-node powers are constant for each scale, carbon
intensity is constant, and there are no waits or restart/checkpoint overheads.
Every execution of the same work, including one that switches scales, has a
time--carbon pair in the convex hull of the fixed-scale pairs.
\end{proposition}
\noindent\emph{Proof.}
Let $x_n$ be the fraction of updates executed at scale $n$, so
$x_n\geq0$ and $\sum_nx_n=1$. If $T_n=U_\star/q_n$ and $C_n$ are the
fixed-scale totals, additivity gives
$(T,C)=\sum_nx_n(T_n,C_n)$.\hfill$\square$

This is a statement about mean time--carbon geometry, not a per-job deadline
guarantee for randomized fixed policies. A preplanned sequence can complete
individual runs at an intermediate time--carbon point even when a per-run
fixed mixture violates a deadline. Thus beating a feasible fixed mixture
alone does not prove that queue feedback is useful. Under real queueing, a request changes
both its current wait and the time and scheduler state at which the next
request arrives; the linear decomposition need no longer hold. Likewise,
time-varying carbon intensity makes an update's carbon depend on its execution
window. These effects create an opportunity, not a guarantee of an adaptive
gain. Our main comparison retains all fixed scales, a deployable mixture,
and a plan-once controller; action-switch frequency alone is not evidence
of improvement.

\paragraph{What the slider can establish.}
Let $V_{\cal P}(D)$ be the infimum of the endpoint objective in
Equation~\ref{eq:objective} over feasible policies in class $\cal P$. These classes describe admissible
controllers with history, not the expressiveness of a particular network.
Fixed policies and their per-run mixtures are contained in preplanned
sequences, which are contained in feedback policies. Therefore
$V_{\rm feedback}(D)\leq V_{\rm preplanned}(D)\leq V_{\rm fixed}(D)$.
Each ideal envelope is nonincreasing in $D$, since relaxing the budget retains
every formerly feasible policy. This is an opportunity bound, not a guarantee
that finite PPO training finds the envelope. Strict improvement requires
useful timing or admission variation; the null case above rules out claiming
it from scaling inefficiency alone. A sequence that beats both fixed points
still needs a preplanned comparator to establish the value of feedback.

\section{Carbon Exposure with Unmeasured Power}
\label{sec:accounting}

Fix one machine--workload pair; its indices are suppressed below. The mean
complete-node power assigned to scale $n$ is
\begin{equation}
 \widehat P_n=P_0\kappa s_n,\qquad
 s_n=\rho+(1-\rho)\eta_n,\quad
 \eta_n=\frac{W_4}{W_n},
 \label{eq:power}
\end{equation}
where $W_n=nU_{\rm profile}/q_n$, $\kappa>0$, $s_4=1$, and
$0\leq\rho\leq1$. The reference $P_0=1$\,kW fixes a normalization unit; $\kappa$ absorbs the
unknown workload-specific node-power level. Neither is a reported A800 server
measurement. We report within-panel carbon ratios and node-hours, rather than
absolute carbon or specification-based node-power estimates.
The workload coefficient $\kappa$ sets the four-node power level. The scale
factor interpolates between constant per-node power ($\rho=1$) and a
constant-useful-work dynamic-energy limit ($\rho=0$, without overhead).
$\eta_n$ is a scaling-efficiency ratio, not measured utilization, and $\rho$
is not the CPU fraction of node power. We use this family as a sensitivity
model; its parameter interval is a declared scenario, not a confidence interval.
At constant CI and without overhead, Equation~\ref{eq:power} gives total energy
$P_0\kappa[\rho W_\pi+(1-\rho)W_4^\star]$, where $W_\pi$ is the
execution's node-hours and $W_4^\star$ is the cost of doing all target updates
at four nodes. For $\rho>0$, changing this parameter cannot reverse an energy
ranking by node-hours. It changes the magnitude of savings in this limit,
not which execution uses less energy.

For an execution log $\pi$, let $\mathcal I_{\pi,n}$ contain all intervals
allocated to the target at scale $n$, including restart and checkpoint. With
realized average grid intensity $c(t)$, define the scale-resolved exposure
\begin{equation}
 L_{\pi,n}=n\int_{\mathcal I_{\pi,n}}c(t)\,dt,
 \qquad
 C_\pi(\rho)=P_0\kappa\sum_n s_n L_{\pi,n}.
 \label{eq:exposure}
\end{equation}
Power is in kW and time in hours. For the primary EIA scenario, intensity is
in kgCO$_2$/kWh; absolute $C_\pi$ remains unknown without $\kappa$.
Non-CO$_2$ greenhouse gases are
excluded. Replacing $c(t)$ with 1 yields IT energy; summing
allocated node time yields node-hours without a power assumption. Applying
the same mean $\widehat P_n$ to all allocated phases is an explicit
approximation. Phase durations are retained in the log so its effect can be
examined without claiming phase-specific power measurements.

Writing $A_\pi=\sum_nL_{\pi,n}$ and
$B_\pi=\sum_n\eta_nL_{\pi,n}$ gives
\begin{equation}
 C_\pi(\rho)=P_0\kappa
              [B_\pi+\rho(A_\pi-B_\pi)].
 \label{eq:affine}
\end{equation}
\begin{proposition}[Endpoint comparison]
\label{prop:endpoints}
For fixed policies and execution logs, $C_\pi(\rho)\leq C_b(\rho)$ for every
$\rho\in[\rho_-,\rho_+]$ if and only if this inequality holds at both
endpoints. The same statement holds for expected carbon under a fixed
distribution of arrivals and policy randomness.
\end{proposition}
\noindent\emph{Proof.}
The difference is affine in $\rho$ by Equation~\ref{eq:affine}; its maximum
on a closed interval occurs at an endpoint. Expectation preserves the affine
form.\hfill$\square$

The common positive factor $P_0\kappa$ cancels from within-panel
ratios and rankings. No workload coefficient sweep is needed for these
comparisons, although unknown coefficients preclude pooling absolute carbon
across workloads. For a positive affine reference $C_{\rm ref}(\rho)$, the
ratio $\mathbb E[C_\pi(\rho)]/C_{\rm ref}(\rho)$ is linear-fractional and
has a derivative of constant sign (or zero). Its worst value also occurs at
an endpoint. Policies must remain fixed while computing this result: retraining
at each $\rho$ would evaluate different policies, not certify one policy.

For powers outside this family, Equation~\ref{eq:exposure} remains useful.
If each scale exposure is no greater than a baseline's, dominance holds for
all positive scale-power coefficients. Otherwise a reversal is possible,
and the result must retain a power-model qualification. A more informative
local test allows $s_n=s_n^0(1+\delta_n)$ for $n\ne4$, with
$|\delta_n|\leq d<1$ and $s_4=1$. Writing
$\Delta L_n=L_{\pi,n}-L_{b,n}$, the worst normalized difference is exactly
\begin{equation}
 \max_{\delta}\frac{C_\pi-C_b}{P_0\kappa}
 =\sum_n s_n^0\Delta L_n
       +d\sum_{n\ne4}s_n^0|\Delta L_n|.
 \label{eq:box}
\end{equation}
This tests independent scale-power errors using the same logs, without
adding a power-model layer. Queueing draws no
allocated energy at this accounting boundary. We exclude embodied emissions,
idle and background-job power, and shared facility overhead. The estimated
quantity is job-attributed operational carbon, not avoided grid emissions.

\section{\name{} Controller}
\label{sec:method}

\subsection{A Budget and a Causal Observation}
A slider $s\in[0,1]$ selects a completion budget
$D(s)=T_{\rm ref}[\beta_{\min}+s(\beta_{\max}-\beta_{\min})]$.
Its left end requests a tighter finish; its right end permits more time to
reduce carbon. The supported ticks are frozen before evaluation; intermediate
positions require validation, not curve interpolation. The miss tolerance
$\epsilon$ and power interval are experiment settings, not additional sliders.
The chosen $D$ is fixed for the run. At boundary $k$, actual remaining time is
$b_k=D-(t_k-t_0)$, which may be negative. The policy observes
\begin{equation}
 o_k=\left[U_k/U_\star,\ b_k/T_{\rm ref},\ D/T_{\rm ref},\ z_k,\
                    g_k,\{\phi_{k,n}:n\in\mathcal A\}\right],
 \label{eq:state}
\end{equation}
where $T_{\rm ref}>0$ is a frozen training-set time reference and $z_k$
summarizes visible queue history and calendar time. Each $\phi_{k,n}$ contains
node fraction, normalized rate $q_n$, planned work $v_{k,n}/U_\star$, allocation
duration $d_{k,n}$, requested limit, and efficiency $\eta_n$. These descriptors
include the shorter final chunk. The vector $g_k$ contains 28 consecutive
six-hour averages of a causal CI forecast, starting at $t_k$, normalized by a
positive training reference. A rolling past hour-of-week forecast supplies
this one-week horizon; the controller is not given a predicted start time.

\paragraph{Decisions without an admission predictor.}
The main actor and critic consume neither predicted waits nor exposure
integrated over a predicted execution window. They learn continuation costs
from complete replay outcomes in one end-to-end training procedure. No
supervised admission model precedes policy training. After each completed allocation, actual elapsed
time updates $b_k$ and completed work updates $U_k$ before the next decision.
Only physical work feasibility masks actions; a predicted deadline cannot
exclude a request. This removes an explicit wait-calibration prerequisite,
but the learned policy still depends on relationships present in its training
environment. It neither eliminates uncertainty nor establishes transfer to
unseen scheduling rules. Feedback cannot recover a wait already incurred.

The queue history uses counts, node demand, requested limits, elapsed waits,
occupied nodes, and elapsed running time at fixed lags. It never uses a running
job's eventual completion time. Missing fields have explicit masks; site
visibility must be checked rather than inferred from trace completeness.
The observation is partially observed: hidden Slurm state need not be
recoverable or Markov given $o_k$.

\subsection{Learning across Power Scenarios}
For endpoint $e\in\{-,+\}$ let
$J_e(\pi)=\mathbb E[C_\pi(\rho_e)]/C_{{\rm ref},e}$, where
$C_{{\rm ref},e}>0$ is the training-set mean for Fixed-4 under that endpoint.
Let $J_V(\pi)=\Pr_\pi(T>D)$. For each configured budget, the objective is
\begin{equation}
 \min_\pi\max_{e\in\{-,+\}}J_e(\pi)
       \quad\text{subject to}\quad J_V(\pi)\leq\epsilon .
 \label{eq:objective}
\end{equation}
The reference and all normalizers are frozen before validation. By
Proposition~\ref{prop:endpoints} and the ratio property following it, the maximum
also covers the interior of the chosen $\rho$ interval for this same policy.
This is robustness to power-accounting uncertainty; it does not imply
robustness to arbitrary queue distributions. The chance constraint is a
population target, not a guarantee for an individual job. If Slurm withholds
all feasible allocations past $D$, every policy under this interface misses;
no user-level controller can guarantee completion against unrestricted waits.

We optimize Equation~\ref{eq:objective} through
\begin{equation}
 \min_\theta\max_{\substack{p\in\Delta_2\\\lambda\geq0}}
       \sum_e p_eJ_e(\pi_\theta)
                    +\lambda[J_V(\pi_\theta)-\epsilon],
 \label{eq:lagrangian}
\end{equation}
where $\Delta_2$ is the probability simplex. At training iteration $j$,
freeze $p,\lambda$ for rollout and PPO updates; then update
\begin{align}
 p_e&\leftarrow
   \frac{p_e\exp(\eta_p\widehat J_e)}
        {\sum_f p_f\exp(\eta_p\widehat J_f)},\nonumber\\
 \lambda&\leftarrow
       [\lambda+\eta_\lambda(\widehat J_V-\epsilon)]_+ .
 \label{eq:dual}
\end{align}
Budget-specific dual variables prevent an easy budget from hiding violations
at a tight one. One actor is conditioned on the budget; models are trained
separately per machine--workload panel, without a cross-workload transfer claim.

A shared two-layer, 128-unit ReLU action encoder and mean-pooled context
produce four masked logits; the value network uses the same width.
Three value heads estimate the remaining carbon cost at the two endpoints
and the probability of an eventual deadline miss. For chunk $k$, the actor
reward is
\begin{equation}
 r_k=-\sum_e p_e\frac{C_k(\rho_e)}{C_{{\rm ref},e}}
        -\lambda\,\mathbf1\{k=K-1,\ T>D\}.
 \label{eq:reward}
\end{equation}
The carbon increment includes the entire allocation. All returns use
$\gamma=1$ and complete episodes. Their sum is exactly the negative
Lagrangian cost up to the constant $\lambda\epsilon$, regardless of how many
chunks the run needs. A fixed per-chunk discount would instead favor moving
cost to later chunks. Value heads fit undiscounted Monte Carlo returns;
the weighted cost advantage gives the PPO update. A missed deadline does not
terminate training work or discard the job's later carbon.

This construction reuses each execution to train both carbon heads and the
miss head; it does not require a simulator run for every power parameter.
The maximization is over expected episode cost, not a separate worst endpoint
chosen at every chunk. PPO and the dual updates are numerical approximations,
so validation and test must evaluate the attained carbon--deadline tradeoff.

\subsection{Selection and Submission}
Before evaluation, we freeze an empirical selection rule: among complete
validation cohorts with observed miss rate at most $\epsilon$, select the
checkpoint minimizing the endpoint objective. Apply the same rule to baseline
selection. If none qualifies, retain the complete candidate with the smallest
miss rate as a labeled fallback. Empirical target attainment is distinct from
a confidence bound supporting population feasibility; an inconclusive bound
does not prove infeasibility. Test reports every configured budget and seed.
The deployed action
rule remains the same categorical sampling rule used for evaluation; changing
to argmax requires a separate validation because it changes the policy.

The execution interface needs a profile, budget, visible queue snapshots, and
a causal CI feed. It writes the selected node count and time limit into the
submission and invokes \texttt{sbatch}~\cite{slurmsbatch}. The training process
must save its checkpoint before allocation expiry; a wrapper cannot checkpoint
a process after Slurm has killed it. Resubmission occurs only after a completed
checkpoint and terminal job status are confirmed. A job-ID log prevents duplicate
submission after wrapper restart. Missing inputs select a preconfigured fixed
policy, whose outcomes remain part of evaluation. This describes the intended
interface; live deployment is not an established result of this draft.

\section{Evaluation}
\label{sec:evaluation}

\noindent\textbf{Evidence status.} The following four packages are reserved
experiment designs and result slots. They specify the minimum evidence needed
for the paper's claims; they do not report completed experiments. The historical data were previously used in a train/validation split. We
therefore use a retrospective chronological benchmark, disclosing prior use
and fixing this revision's train/validation/test intervals before execution.
Fit on training and select on validation; do not revise methods using the
new test outcomes. This protocol does not create a previously unseen holdout.

\subsection{E1: Are the Work and Contention Scenarios Sound?}
\slot{E1: inputs and scenario contract}{
Use Figure~12 of AMSP's March 2024 revision~\cite{chen2024amsp}: LLaMA-7B/13B/30B,
BF16, 4096-token sequences, and fixed global batch. Its 32/128/512/1024-GPU
points map to 4/16/64/128 eight-A800 nodes. Archive PDF hash, vector-bar
extraction, units, and the simulation configuration. Audit input workloads,
state initialization, work accounting, and independent wait probes.
Scenario statistics and validation outcomes: \pending{E1}.}

\paragraph{What the traces represent.}
A background job supplies a tuple of submission time, node demand, requested
wall time, and occupied duration. These tuples define rigid, homogeneous-node
workloads with no modeled network or I/O interference. Each source's time
order and joint width--duration structure remain intact; traces are not pooled.
The nominal cluster has 128 nodes, matching the profile testbed's size. This
is a declared experimental capacity, not an estimate of either source cluster.
Old durations specify scenario demands; they do not predict the old programs'
runtimes on A800s. Queue state and waits are regenerated under the declared
scheduler. Conditional on identical tuples, initial state, and scheduler,
changing hardware labels cannot change admission. That conditional property
neither makes workloads hardware-independent nor establishes universal transfer.

\paragraph{Replay and validation.}
The main scheduler is a node-level conservative-backfill approximation, not
Slurm itself. Rebuild background state from an empty state at the declared
trace-prefix origin, excluding at least 28 days before scored arrivals and
checking prefix sensitivity in development. Do not transplant historical starts
or running allocations into the constructed cluster. Methods branch from the
same background state; their requests then affect subsequent admissions.
State continues across all chunks. Original observed waits are not ground
truth for this constructed scenario. Verify scheduling mechanics on
independently solvable cases and evaluate the predictor on held-out simulator
probes; source-calendar tests preserve chronology and report censoring.

\begin{table}[t]
\centering\small
\begin{tabular}{@{}p{0.42\linewidth}p{0.50\linewidth}@{}}
\toprule
Input / check & Reserved evidence \\
\midrule
Trace demand and burstiness & Pending: counts, widths, durations \\
Scenario contention & Pending: occupancy and wait tails \\
Published throughput input & AMSP Fig.\ 12; extraction archived \\
State and work conservation & Pending: scheduling/accounting checks \\
Prediction vs.\ probe wait & Pending: errors and coverage \\
Unmeasured execution costs & Declared overhead scenarios \\
\bottomrule
\end{tabular}
\caption{Input and validity results placeholder. Report both source workloads
separately, including time zones, cleaning, clipping, and calendar blocks.
Probe labels come from the simulator; they are not observed alternative
requests on a production cluster.}
\label{tab:validity}
\end{table}

\paragraph{Execution inputs.}
Figure~12 reports tokens per GPU per second, $z_n$; convert to updates per
hour as $q_n=3600(8n)z_n/(1024\cdot4096)$. The 1024-sequence batch follows
from the source's stated micro-batch and accumulation settings. Extract only
published points; 64- and 256-GPU points are absent. For each model, freeze
$U_\star=\lceil192q_4\rceil$ before scheduling experiments. This chooses a
workload of approximately 192 useful-compute hours on four nodes; overhead
and rounding determine the chunk count, rather than a convergence target. Use 600 seconds of total allocated overhead per chunk, split equally
between initialization/restart and saving, as an explicit scenario. E4 varies
this unmeasured cost. No checkpoint implementation, final-quality equivalence,
or complete-node power measurement is claimed.

\subsection{E2: Is Adaptation Better than Strong Alternatives?}
\slot{E2: primary comparison}{
Evaluate paired initial arrivals for both source workloads and all three published profiles.
Use EIA's Texas ERCOT average consumption-based CO$_2$ intensity as the primary
scenario~\cite{eia2026gridmonitor}. Set $T_{\rm ref}$ to the fastest fixed
training mean TAT and $D_{\rm hi}$ to the training Fixed-4 95th-percentile TAT.
Use $\beta_{\min}=1$, $\beta_{\max}=D_{\rm hi}/T_{\rm ref}$, and
$s\in\{0,1/4,1/2,1\}$. This sets the budget span, not a feasibility guarantee.
Freeze the grid, miss tolerance, power interval, seeds, and selection rule before test access. Compare \name{} with every fixed scale, a validation-fitted
fixed mixture, plan-once control, and rollout planning. Main figure, carbon ratios,
miss rates, TAT, and node-hours: \pending{E2}.}

\paragraph{Baselines.}
Fixed-4/16/64/128 retain the same scale across all chunks and use the identical
work, checkpoint, and final-request rules. \emph{Best-Fixed} is selected on
validation for each budget. \emph{Fixed-Mix} samples one scale at the start of
a run and keeps it: its four mixture weights solve the same endpoint objective
and miss constraint using validation statistics. Mixture feasibility depends
on the mixture of miss probabilities, not on an interpolated mean TAT.
All four fixed points remain visible even if dominated.

For the planning baselines only, fit a request-conditioned boosted regressor
for $\log(1+w)$; chronological training residuals supply 32 equal-mass atoms
of $\widehat F_{k,n}$. Its causal exposure estimate is
\begin{equation}
 \widehat L_{k,n}=n\,\mathbb E_{w\sim\widehat F_{k,n}}
 \left[\int_{t_k+w}^{t_k+w+d_{k,n}}\widehat c_{t_k}(s)\,ds\right].
 \label{eq:forecast}
\end{equation}
E1 reports probe errors and coverage to document this approximation; prediction
accuracy is not an acceptance criterion for \name{}.

\emph{Rollout-MPC} enumerates two upcoming scale choices and a constant-scale
tail to completion, uses the auxiliary wait model and causal CI forecast, and executes
only the first action before replanning. There are at most $4^3$ candidate
plans. Future queue observations are unavailable, so the rollout holds the
current queue and wait-model calendar features fixed, updates request length
and the forecast integration window, and samples conditional waits. Its carbon
estimate adds past endpoint costs using released CI, retaining each chunk's
submission forecast for unreleased intervals, to the planned continuation. It selects the
lowest worst-endpoint normalized carbon plan meeting the estimated miss tolerance, or the
plan with the smallest miss estimate if none qualifies. This is a strong
causal baseline with an explicit approximation, not a clairvoyant scheduler.
\emph{Plan-once} runs the identical planner only at the first submission and
executes its selected scale sequence and constant tail without replanning.
It receives the initial queue and forecast, allowing a planned within-run
mix while removing feedback from later admissions. Both planning variants
may tune their internal risk threshold on validation; test compares them
under the same actual miss tolerance as every other method.

\begin{figure*}[t]
\centering
\figslot{1.65in}{E2: Two source workloads $\times$ three published AMSP profiles.\\
All fixed scales; Fixed-Mix; Plan-once; Rollout-MPC; \name{}.\\
Mean TAT versus modeled carbon at a declared $\rho$; mark observed miss violations and label slider ticks.\\
Companion labels: worst-endpoint carbon ratio and deadline-miss rate.}
\caption{Primary results placeholder. Every point completes the same prescribed
work from paired initial arrivals. Report hours, modeled carbon normalized
within each contention--profile panel, 95\% paired block intervals, the power
scenario, and the frozen budget grid. Mark observed miss-target violations
explicitly. Fixed-Mix is an executable per-run randomized policy; connecting
other points does not imply an evaluated intermediate policy. Report every
slider tick; do not smooth reversals or select settings after test.}
\label{fig:main}
\end{figure*}

\paragraph{Paired analysis and selection.}
An episode is one initial arrival and complete run. Resample contiguous calendar
blocks jointly across methods, with RL seeds as another replication level.
Blocks must cover the history window and development dependence. Training fixes
budgets and normalizers; validation fixes operating points. The primary
comparator is the strongest validation-selected non-RL method meeting the
empirical miss target at the same budget. Fixed-128 is not the automatic
denominator. Confidence bounds are reported separately from this frozen
empirical selection rule.

Report the endpoint ratio of \emph{mean} carbon costs, mean and p95 TAT,
miss rate, and allocated node-hours per panel; macro-averages use dimensionless
panel ratios. Carbon includes all completed runs, including deadline misses.
Episodes crossing a split must be purged or assigned entirely to one split,
and any run not observed to completion is reported as censored, not silently
dropped. Texas CI must have a documented source, release time, time zone,
resolution, and missing-data rule. If CI and queue years do not overlap,
their alignment is a labeled counterfactual pairing.

\paragraph{Result slot and inference.}
\pending{E2: effect sizes, intervals, eligible budgets, and failed panels}.
An observed budgeted carbon reduction requires both frozen methods to meet
the miss target on complete validation and test cohorts. This describes
observed tradeoffs, not population feasibility. A stronger statistical claim
also requires joint endpoint cost intervals and one-sided miss bounds to pass
the declared criteria. Report inconclusive bounds, including for zero misses;
nearby arrivals do not create independent calendar blocks. Gains over
Fixed-128 alone establish neither adaptive nor RL value. An inspected former
test period cannot be relabeled as untouched.

\subsection{E3: Does Feedback Beat a Learned Sequence?}
\slot{E3: matched feedback comparison}{
Reuse E2 \name{}, Plan-once, and Rollout-MPC. Train one additional
\emph{Precommitted-RL} policy with the same architecture, budget grid, seeds,
arrival sampling, and episode-interaction budget. Before its first submission,
it samples an entire scale sequence using only initial public queue/CI inputs,
remaining planned work, and $D$ minus planned allocation time. It receives no
later queue observations, CI releases, or realized elapsed time. Evaluate at
one prespecified budget per panel. Results: \pending{E3}.}

\begin{table}[t]
\centering\small
\begin{tabular}{@{}lccc@{}}
\toprule
Method / variant & Carbon & Misses & Node-h \\
\midrule
\name{} & --- & --- & --- \\
Precommitted-RL & --- & --- & --- \\
Rollout-MPC & --- & --- & --- \\
Plan-once & --- & --- & --- \\
\bottomrule
\end{tabular}
\caption{Feedback results placeholder at prespecified budgets. All dashes
are pending E3 results. Use E2 normalization and paired intervals; compare
carbon jointly with misses. Precommitted-RL changes its scale within a run
but commits its choices before seeing subsequent outcomes.}
\label{tab:ablation}
\end{table}

The learned-sequence comparison removes feedback while retaining the capacity
to mix scales and optimize from completed episodes. Its remaining-time input
is planned allocation time only, not a zero-wait assumption in the execution
environment: all actual waits and costs enter training returns and evaluation.
Archive each full plan before executing it. Neither policy needs an admission
predictor. Matching interaction budgets prevents extra replay from explaining
a feedback gain, although it cannot guarantee equal optimization quality.
Plan-once versus MPC separately tests replanning within an explicit planner.
Report single-chunk frequency: those runs offer no later decision. A fixed-scale
win without a learned-sequence win supports mixing, not a feedback advantage;
a win only over MPC could reflect its modeling error. Action frequencies alone
do not establish either claim. Results: \pending{E3: feedback increment}.

\subsection{E4: Does the Conclusion Depend on the Power Model?}
\slot{E4: robustness from existing logs}{
Freeze policies and recompute every E2 execution over the declared $\rho$
interval, including both endpoints. Report pairwise break-even points,
componentwise exposure differences, and the independent scale-error radius
supported by Equation~\ref{eq:box}. Compare the robust policy with a
single-endpoint-trained counterpart at one prespecified budget per panel.
Reuse existing logs for parameter sensitivity; queue and forecast perturbations
require fresh replay. Outcomes: \pending{E4}.}

\begin{figure}[t]
\centering
\figslot{1.05in}{E4: carbon advantage versus $\rho$.\\
Endpoint intervals; zero-gain line; break-even $\rho$.\\
Fixed policies and comparator throughout.\\
Small inset: endpoint-trained versus robust policy.}
\caption{Power-sensitivity results placeholder. The parameter range is an
assumed scenario, not a measured power range. Endpoint comparisons characterize
the interval only under Equation~\ref{eq:power}. Plot uncertainty from paired
execution samples separately from power-scenario variation.}
\label{fig:power}
\end{figure}

Figure~\ref{fig:power} reserves the sensitivity result.
\pending{E4: interval-wide advantage or the exact region of reversal}.
For $\Delta A=A_\pi-A_b$ and $\Delta B=B_\pi-B_b$, a possible break-even point
is $\rho^\star=-\Delta B/(\Delta A-\Delta B)$ when the denominator is nonzero.
A point outside the declared interval is not an in-range reversal.
Simultaneous one-sided endpoint intervals, obtained from the paired block
resamples, support an interval-wide statistical statement; two point estimates
alone do not. A nominal gain that reverses near the selected power assumption
must be reported as model-sensitive.

\paragraph{Does the simulated environment determine the gain?}
At one prespecified budget and the 7B profile, reuse frozen full-policy and
Rollout-MPC models in each source workload. Independently change (i) background
widths to $\min(2n,128)$, reporting the capped fraction, (ii) scheduling to
strict FCFS, and (iii) total chunk overhead to 60 or 1800 seconds. Preserve
physical budgets, workload, CI calendars, and all fitted parameters. These
four additional settings require new replay, not new training or hardware
measurements. Width and scheduler changes test learned queue relationships;
E1 prediction accuracy alone cannot establish this robustness. Report carbon and misses together: a gain that disappears or
violates the budget limits the claim. They test the selected contention and
overhead scenarios, not every possible cluster. Results: \pending{E4: environment sensitivity}.

A compact reuse analysis reports E2 outcomes by held-out calendar block and
CI variability. A constant-CI rescore of the same logs removes timing-related
carbon differences for those executions; it does not simulate policies reacting
to constant CI. No cross-region, multi-user-adoption, or facility-wide
experiment is needed for the scoped claim.

\section{Limitations and Conclusion}

\paragraph{Limitations.}
Evidence is limited to declared demand and throughput scenarios, including
synthetic GPT-2 efficiency curves. Fixed historical arrivals omit user responses
to changed waits. Unmeasured scheduler details, placement interference,
checkpoint correctness and costs limit hardware transfer; overhead sweeps do
not resolve them. Changed queues or CI forecasts can invalidate learned
decisions. Completion budgets and empirical misses are not production SLAs.

Complete-node power is unmeasured. Certificates require the declared,
time-invariant scale coefficients; common factors cancel only within panels.
Job attribution does not establish facility-wide, embodied-carbon or adoption
impacts. Include training and profiling costs before claiming amortized benefit.

\paragraph{Conclusion.}
\name{} learns checkpoint-boundary resource requests from complete outcomes;
an exposure decomposition makes unknown node power explicit.
\pending{E2--E4: baseline-relative results, feasibility, and supported power interval}.
Gains over fixed scales establish an opportunity. Feedback attribution also
requires gains over a learned precommitted sequence; RL claims require
comparison with explicit planning.

\bibliography{reference,scaledown_refs}
\end{document}
